% Compose the independently proved coding and Pell reductions.
\section{A certificate of 99 operations}\label{sec:composed-99}

The reductions in Sections~\ref{sec:half-center} and
\ref{sec:doubled-index} can be applied together. Their arithmetic
changes occupy disjoint parts of the certificate. The composition
also preserves the order of the mathematical proof: the new packing
first supplies the bounds for the new Pell block; that block then
recovers the exponent and binary radix needed by the new digit tests.

\begin{theorem}\label{thm:best99}
For every recursively enumerable subset of the positive integers,
there is an admissible fixed index $(V,H,T_L)$ such that its membership
relation is represented by a system with 34 positive unknowns and
22 equations admitting a straight-line certificate of 99 operations:
55 multiplications and 44 additions. Here $T_L=\psi_4(L)$ for the
underlying exponent of the index construction. The supplied parameters
are $(x,V,H,T_L)$, with three fixed index parameters, and fixed numerals
are free.
\end{theorem}

\begin{proof}
Construct the coefficient index using the entire paired-row
construction of Section~\ref{sec:half-center}. In particular, the
special target $\delta^2$ comes first; multiplication, addition,
copy, equality, and zero rows are paired with their negatives; each
unit gate uses the two unpaired rows
$X\delta-\delta^2$ and $X\delta-X^2$; and the guard is
$u\delta-x^2$. Every ordinary target is $2F+\delta^2$.

If there are $m$ positive-weight coordinates and $s$ targets, use
\[
\begin{gathered}
 v_i=3^i\quad(0\le i<m),\qquad M=3^{m-1},\qquad d_0=4M+1,\\
 t_j=(s+1)d_0+2M+jd_0\quad(0\le j<s),\\
 t_*=2sd_0+2M,\qquad K_0=t_*+1,\qquad c_*=m+s+1.
\end{gathered}
\]
Choose a power of two $L>3K_0+2$, form the corresponding
$\ell_0,e_0$, and set
\[
 V=\ell_0(4)+e_0(4)4^L,\qquad
 H>\max\{2\,4^{2L+1},\ 4^{t_*+3}\,4c_*^2,\ 3L,\ 16\},
\]
with $H$ a power of two. Finally supply $T_L=\psi_4(L)$ instead
of $L$. Thus $V,H$ come from the new paired layout throughout;
they are not reused from an earlier unpaired index. These choices
are independent of the queried input $x$.

Keep the coding equations and packing
\eqref{eq:half-center-block}--\eqref{eq:half-center-packing}, with
\[
 \Omega=2\lambda-e>0,\qquad
 \sigma=\theta\lambda q-\Omega\mathcal C^2>0,
 \qquad \theta=B-4,
\]
and apply both Pell changes
\eqref{eq:doubled-auxiliary} and \eqref{eq:doubled-fixed-index}.
All other common-witness Pell equations remain unchanged.

Before either code is decoded, the positive bound $S_2+\alpha=q^2$
gives $e<q$ and $\ell<q^2$. The geometric equation gives
$B\le q^2$, and the index bound gives $q>4b\ge8$ and
$3L\le B\le n$. Positivity of $\Omega$ and $\sigma$ gives
\[
 \mathcal C^2<\theta\lambda q<3q^3<q^4.
\]
The preliminary borrow argument of Section~\ref{sec:half-center}
therefore proves
\begin{equation}\label{eq:composed99-preliminary}
 0<S<n,\qquad 0<T+1\le n,\qquad
 n\le r\le2n^3-2n^2<2n^3.
\end{equation}
No coefficient target has been tested at this stage.

Put $U=wn^2$, $Y_P=sn^2$, $a=Y_P(U+1)$, $A=a+4$, and $J=2r+1$.
Then $U,Y_P\ge n^2$, $UY_P\ge n^4>r+1$, and the first-index
congruence and interval give $c>J$. Also $A>n^4>J$.
Thus every hypothesis of Lemma~\ref{lem:doubled-main-index} holds:
$A>1$, $1<J<c$, and $J$ is odd. That lemma proves
\[
 c=\psi_A(J),\qquad d=\chi_A(J).
\]
In particular, its signed comparison does not depend on either
coefficient decoding or an already known value of $q$.

The unchanged first Pell index and lower ratio argument now give,
with $\xi=(U+1)^{2r}/U^r$,
\[
 k=\psi_{2UY_P^2+1}(r+1),\qquad
 \xi<c/k<Y_P+1,\qquad
 Y_P\ge U^r,\qquad a>U^{r+1}.
\]
The first exponential congruence yields $U=4^J$. Therefore $n$ and
$q$ are powers of two, and $a>4^{3J}$.
The positive gap and second norm give
$\kappa=\psi_A(t)$ for some $0<t<J$. Admissibility gives
$L<J$. Lemma~\ref{lem:doubled-fixed-index} applies to
$\kappa=T_L+\Delta a$: reduction modulo $a$ gives
$\psi_4(t)\equiv\psi_4(L)$, while both values lie strictly between
zero and $a$. It follows that $t=L$.

The remaining exponent comparison uses
\[
 B^{3L}\le B^B\le n^n\le U^r<A,
 \qquad q^3\le n^n\le U^r<A,
\]
so it yields $q=B^L$, and hence powers of two $B$ and $b$.
The unchanged upper ratio and binomial arguments give
\begin{equation}\label{eq:composed99-binomial}
 n^2\mid\binom{2r}{r}.
\end{equation}
These conclusions precede all internal circuit tests.

Now $\lambda>q$ and $e<q$, so $\Omega>\lambda$ and
$\mathcal C^2<\theta q<Bq<q^2$. Thus $g<\mathcal C<q$.
All post-Pell hypotheses of Section~\ref{sec:half-center} have been
restored. Its first two masks remove the preliminary borrow and
identify the packed coefficient polynomial, up to the quotient alias.
The first mask bounds every coordinate by $b$. Positivity of $x,g$
gives $\mathcal C(1)\ge2$ before $\delta$ is identified. Accordingly,
with $A_{\mathcal C}=\mathcal C(1)^2$, the modified high-window bound
$2A_{\mathcal C}+4\le4A_{\mathcal C}$ and the displayed bound on $H$
remove the alias. The two codes are canonical.

The first special target has incoming carry zero and forces
$\delta\in\{0,1\}$. The guard excludes zero. At $\delta=1$, each
ordinary row tests $F\in\{0,1\}$; its paired negative forces $F=0$,
and the two unit rows force the unit coordinate to equal one.
This recovers the represented system at the actual input $x$.

For the converse, take a solution of that system, its circuit values,
$\delta=1$, $u=x^2$, and zero dummy digits. Choose a power-of-two
$b$ larger than every coordinate and use the canonical codes.
Then $\beta,\theta,\lambda,\alpha$ and the packed-congruence quotient
are positive integers. The half-center necessity proof gives positive
$\Omega,\sigma$ and all three no-carry conditions. Define $r$ by
the packing; \eqref{eq:composed99-preliminary} and
\eqref{eq:composed99-binomial} follow.

Choose $J=2r+1$, $U=4^J$, and $w=U/n^2$. This quotient is a positive
integer since $n$ is a power of two and $n^2\le r^2<4^J$.
Take $Y_P=\lfloor\xi\rfloor$ and $s=Y_P/n^2$, integral by the
binomial congruence, and define $a=Y_P(U+1)$.
The common-witness construction supplies positive
$k,c,d,\tau,h,\eta,\zeta,\gamma,\rho$.
Choose
\[
 \kappa=\psi_A(L),\quad \mu=\chi_A(L),\quad
 \phi=c-\kappa,\quad
 \Delta=\frac{\psi_A(L)-\psi_4(L)}{a}.
\]
Here $L<J$ makes $\phi>0$, and $A=a+4>4$, $L\ge2$ make
$\Delta$ a positive integer.

Finally choose the auxiliary witnesses exactly as in
Lemma~\ref{lem:doubled-main-index}: set
\[
\begin{gathered}
 m_0=2cJ,\qquad f=\chi_A(m_0),\qquad
 i=\psi_A(m_0)/c^2,\\
 K=(A^2-1)(f^2-1),\qquad G=2K-1,\qquad
 I=\chi_G(J),\qquad H_0=\psi_G(J),\\
 o=\frac{I+2d^2-1}{2f},\qquad j=\frac{H_0-J}{c}.
\end{gathered}
\]
The proved divisibility, oddness of $f$, and Pell congruences make
all four auxiliary unknowns positive integers and establish the
doubled norm. This last choice changes no coding, interval, gap,
or exponential witness. Thus every required positive unknown exists.

For the arithmetic count, the half-center substitution saves one
multiplication from Theorem~\ref{thm:best101}; the doubled-index
substitution saves one addition. They change disjoint source rows and
disjoint instruction groups. The exact checker
\texttt{round26\_1980\_composed\_certificate.py} verifies that both
orders give the same complete schedule and source system, together
with all primitive instructions, 22 equality tests, and three
triangular residual corrections. The resulting count is
$(56-1)+(45-1)=55+44=99$.
\end{proof}
